\documentclass[aspectratio=169,10pt]{beamer}
\usepackage[utf8]{inputenc}
\usepackage[T1]{fontenc}
\usepackage[spanish,es-nodecimaldot,es-noquoting]{babel}
\usepackage{amsmath,amssymb}
\usepackage{booktabs}
\usepackage{listings}
\usepackage{tikz}
\usetikzlibrary{babel,arrows.meta,positioning,calc}
\usepackage{pgfplots}
\pgfplotsset{compat=1.17}

\usetheme{default}
\usecolortheme{default}
\setbeamertemplate{navigation symbols}{}
\setbeamertemplate{footline}[frame number]
\definecolor{acento}{RGB}{0,98,128}
\definecolor{acento2}{RGB}{214,110,30}
\setbeamercolor{frametitle}{fg=acento}
\setbeamercolor{title}{fg=acento}
\setbeamercolor{structure}{fg=acento}
\setbeamerfont{frametitle}{series=\bfseries}

\lstset{language=C++,basicstyle=\ttfamily\scriptsize,keywordstyle=\color{acento}\bfseries,
  commentstyle=\color{gray},showstringspaces=false,frame=single,rulecolor=\color{gray!40},
  columns=fullflexible,literate={á}{{\'a}}1 {é}{{\'e}}1 {í}{{\'i}}1 {ó}{{\'o}}1 {ú}{{\'u}}1 {ñ}{{\~n}}1}

\newcommand{\ok}{\textcolor{green!50!black}{\checkmark}}

\title{Estado de Avance\\\large Bus Driver Scheduling Problem con HC-BI + Restart}
\author{Vicente Luongo \quad ROL 202073637-5}
\institute{INF295 Inteligencia Artificial -- UTFSM}
\date{2026-2}

\begin{document}

\begin{frame}
\titlepage
\end{frame}

%---------------------------------------------------------------
\begin{frame}{Problema y técnica asignada}
\begin{columns}[T]
\column{0.55\textwidth}
\textbf{BDSP} (variante de Constantino et al.~[1]):
\begin{itemize}
\item Cubrir cada tramo con \textbf{exactamente un} conductor.
\item Respetar las reglas laborales: stretch $\le 360$, descanso $\ge 90$, umbral de descanso pagado de 300, sobretiempo $\le 120$, spread $\le 780$.
\item \textbf{Minimizar el tiempo pagado total}, que incluye el mínimo de 432, el recargo nocturno, el recargo por sobretiempo y la penalización por cambio de línea.
\end{itemize}
\column{0.42\textwidth}
\textbf{Técnica: HC-BI + Restart}
\begin{itemize}
\item Representación \emph{ad-hoc}.
\item Inicialización \textbf{aleatoria}.
\item Movimientos que permitan visitar todo el espacio de búsqueda.
\item \emph{Best Improvement}: se explora el vecindario completo y se toma el mejor vecino.
\item Reinicio al llegar a un óptimo local.
\end{itemize}
\end{columns}
\end{frame}

%---------------------------------------------------------------
\begin{frame}[fragile]{Estado actual de la implementación (C++17, Linux)}
\begin{columns}[T]
\column{0.5\textwidth}
\begin{tabular}{ll}
\toprule
Componente & Estado\\
\midrule
Lectura de instancia & \ok\\
Función de evaluación & \ok\\
Solución inicial aleatoria & \ok\\
Escritura de la salida & \ok\\
HC-BI (vecindarios N1, N2) & \ok{} preliminar\\
Reinicios y control de tiempo & \ok\\
Experimentos y ajuste & pendiente\\
\bottomrule
\end{tabular}

\vspace{1em}
Probado con las 10 instancias (130 a 2313 tramos) y con el ejemplo de 6 tramos del curso, que da FO $=864$, la misma cifra del informe.
\column{0.45\textwidth}
\begin{lstlisting}[language=bash]
codigo/
  Makefile    README.md
  src/instancia.{h,cpp}   lectura
  src/evaluacion.{h,cpp}  FO
  src/solucion.{h,cpp}    repr. +
                          inicial +
                          salida
  src/hillclimbing.{h,cpp}
  src/main.cpp
  instancias/

$ make
$ ./bdsp inst.txt <semilla> <seg> [-v]
\end{lstlisting}
\end{columns}
\end{frame}

%---------------------------------------------------------------
\begin{frame}[fragile]{Lectura de la instancia y almacenamiento interno}
\begin{columns}[T]
\column{0.38\textwidth}
\textbf{Entrada:}
\begin{lstlisting}[language=bash]
130
"011 ","G","T",6,5,35,12
"011 ","T","T",6,40,40,12
...
360  90  300  432  120
780  60/52  20  50  1
\end{lstlisting}
\begin{itemize}\small
\item Se toleran CRLF, comillas y espacios.
\item $f_N$ se lee como fracción ($60/52$).
\item Si el formato es inválido, el programa da un mensaje de error.
\end{itemize}
\column{0.58\textwidth}
\begin{lstlisting}
struct Tramo {
  std::string linea; int lineaId; // línea -> entero
  int origen, destino;            // lugar -> entero
  int inicio, fin, duracion;      // minutos desde 00:00
  int bus;
  int nocturno;   // nu_j precalculado (22:00-05:00)
};
\end{lstlisting}
\begin{itemize}\small
\item Hora y minuto se convierten a minutos: $s_j = 60h+m$, $e_j=s_j+d_j$. Las horas $\ge 24$ se aceptan (cruzan la medianoche).
\item $\nu_j=|[s_j,e_j)\cap\mathcal N|$ con $\mathcal N=\bigcup_d[1440d-120,1440d+300)$: se calcula \textbf{una sola vez} al leer la instancia.
\item Las líneas y los lugares se mapean a enteros, de modo que las comparaciones cuestan $O(1)$.
\end{itemize}
\end{columns}
\end{frame}

%---------------------------------------------------------------
\begin{frame}{Representación}
\begin{columns}[T]
\column{0.5\textwidth}
\begin{itemize}
\item \textbf{Asignación directa tramo $\to$ conductor}, que corresponde a las variables $x_{jk}$ del modelo del informe.
\item Internamente: \texttt{vector<vector<int>> turnos}; cada turno guarda sus tramos \textbf{ordenados por hora de inicio}.
\item El orden lo fijan los horarios, así que la asignación define la solución completa: secuencia, descanso, sobretiempo y pago.
\item Cada turno guarda su evaluación en caché (\texttt{EvalTurno}).
\item La cobertura exacta está garantizada por construcción, porque cada tramo aparece en un solo turno.
\end{itemize}
\column{0.47\textwidth}
\begin{tikzpicture}[x=0.36cm,y=0.75cm,font=\scriptsize]
  \draw[->,gray] (0,-0.3) -- (16,-0.3) node[right] {t};
  \node[left] at (0,2) {$k=1$};
  \node[left] at (0,1) {$k=2$};
  \node[left] at (0,0.2) {};
  \foreach \a/\b/\l in {1/2/5,2/4/6}{\draw[fill=teal!40] (\a,1.8) rectangle (\b,2.2); \node at ({(\a+\b)/2},2) {\l};}
  \foreach \a/\b/\l in {8/9/1,9/11/2,11/13/3,13/14.5/4}{\draw[fill=acento2!40] (\a,0.8) rectangle (\b,1.2); \node at ({(\a+\b)/2},1) {\l};}
  \node[align=left,anchor=north west] at (0,-0.6) {\texttt{turnos = \{\{5,6\},\{1,2,3,4\}\}}\\[2pt]
     \texttt{salida: 1 2 \textbar{} 2 2 \textbar{} 3 2 \textbar{} 4 2 \textbar{} 5 1 \textbar{} 6 1}};
\end{tikzpicture}

\vspace{0.5em}
\small Espacio de búsqueda: a lo más el número de Bell $B_n$ (particiones de los tramos).
\end{columns}
\end{frame}

%---------------------------------------------------------------
\begin{frame}{Manejo de restricciones}
\begin{tabular}{p{0.32\textwidth}p{0.2\textwidth}p{0.4\textwidth}}
\toprule
Restricción & Tipo & Cómo se maneja\\
\midrule
Cobertura exacta & dura & La garantiza la representación.\\
No traslape temporal & dura & No se generan vecinos con traslape (búsqueda binaria en $O(\log L)$).\\
\midrule
Continuidad espacial ($b_i=a_j$) & blanda & 60 min de violación por salto de lugar.\\
Spread $\le S^{\max}$ & blanda & Penalización por minuto de exceso.\\
Stretch $\le T^{\max}$ y descanso $\ge R^{\min}$ & blanda & Se busca una pausa válida; si no hay, se penaliza la menor violación.\\
Sobretiempo $\le O^{\max}$ & blanda & Penalización por minuto de exceso.\\
\bottomrule
\end{tabular}

\vspace{0.8em}
$\text{costo}(k)=\text{pagado}_k + p_L\cdot\text{cambios}_k + \mathbf{1000}\cdot\text{violación}_k$

\small El peso 1000 es mayor que $P^{\min}=432$, así que siempre conviene más reparar un minuto de violación que ahorrar un turno.
Las restricciones blandas permiten pasar por soluciones infactibles al mover tramos: por ejemplo, quitar un tramo intermedio de un turno puede romper su continuidad espacial.
\end{frame}

%---------------------------------------------------------------
\begin{frame}[fragile]{Función de evaluación (por turno, en $O(L)$)}
\begin{columns}[T]
\column{0.47\textwidth}
\begin{align*}
\text{spread}_k &= E_k - S_k\\
\text{desc}_k &= \max_{\text{pausa válida}} \min(g,R^{\text{pag}}) \quad(\text{si spread} > T^{\max})\\
W_k &= \text{spread}_k - \text{desc}_k + (f_N-1)\,\nu_k\\
O_k &= \max(0,\ W_k - P^{\min})\\
\text{pagado}_k &= \max(W_k,P^{\min}) + \beta_O O_k + \beta_N f_N \nu_k
\end{align*}
\small Corresponde a la función objetivo del Modelo~2 del informe. Una pausa es válida si dura $\ge R^{\min}$ y deja ambos stretches $\le T^{\max}$. Entre las válidas se elige la que \textbf{maximiza} la parte no pagada.

\vspace{0.3em}
\textbf{Verificación a mano} (V\_TCCC0130, semilla 3, conductor 1):\\
spread 530, descanso 120, $\nu=20$:\\
$W=413{,}1$; pagado $=432+0{,}2\cdot\frac{60}{52}\cdot20=436{,}6$ \ok
\column{0.5\textwidth}
\begin{lstlisting}
EvalTurno evaluarTurno(const Instancia& I,
                       const vector<int>& turno) {
  ...
  if (r.spread > P.maxStretch) {   // regla 1-2
    for (cada pausa g entre a y a+1) {
      v = max(0, Rmin - g)
        + max(0, (fin_a - S) - Tmax)
        + max(0, (E - ini_a1) - Tmax);
      if (v == 0) mejor = max(mejor, min(g, Rpag));
      else        menorV = min(menorV, v);
    }
    ...
  }
  r.trabajado  = spread - descanso + (fN-1)*noct;
  r.sobretiempo = max(0.0, r.trabajado - Pmin);
  r.pagado = max(r.trabajado, Pmin)
           + bO*r.sobretiempo + bN*fN*noct;
  r.costo = r.pagado + pL*cambios
          + PESO_VIOLACION*r.violacion;
}
\end{lstlisting}
\end{columns}
\end{frame}

%---------------------------------------------------------------
\begin{frame}{Solución inicial aleatoria}
\begin{columns}[T]
\column{0.52\textwidth}
\begin{enumerate}
\item Se baraja el orden de los tramos (\texttt{std::shuffle} con la semilla dada).
\item Para cada tramo $j$, en ese orden:
  \begin{itemize}
  \item se calculan los \emph{candidatos}: los turnos donde $j$ cabe sin traslape \textbf{y} sin violar ninguna regla;
  \item si hay candidatos, se elige uno \textbf{al azar} (uniforme);
  \item si no hay, se abre un conductor nuevo.
  \end{itemize}
\end{enumerate}
\vspace{0.5em}
\begin{itemize}\small
\item Siempre entrega una solución \textbf{factible}, por lo que el HC parte desde la región factible.
\item Es aleatoria y no \emph{greedy}: no elige el mejor candidato, sino uno cualquiera. Así cada reinicio parte de un punto distinto.
\item Costo $O(n\cdot m\cdot L)$: menos de 0,03 s incluso con 2313 tramos.
\end{itemize}
\column{0.45\textwidth}
\begin{tikzpicture}[font=\scriptsize,node distance=0.45cm,
  b/.style={draw,rounded corners,fill=acento!10,minimum width=4.2cm,align=center}]
  \node[b] (a) {barajar tramos};
  \node[b,below=of a] (b) {siguiente tramo $j$};
  \node[b,below=of b] (c) {candidatos = turnos donde\\$j$ cabe y todo es factible};
  \node[b,below=of c] (d) {¿hay candidatos?};
  \node[b,below left=0.5cm and -1.6cm of d,minimum width=1.8cm] (e) {elegir uno\\al azar};
  \node[b,below right=0.5cm and -1.6cm of d,minimum width=1.8cm] (f) {conductor\\nuevo};
  \draw[-Stealth] (a)--(b); \draw[-Stealth] (b)--(c); \draw[-Stealth] (c)--(d);
  \draw[-Stealth] (d)--node[left]{sí}(e); \draw[-Stealth] (d)--node[right]{no}(f);
\end{tikzpicture}
\end{columns}
\end{frame}

%---------------------------------------------------------------
\begin{frame}{Algoritmo: Hill-Climbing Best Improvement + Restart}
\begin{columns}[T]
\column{0.5\textwidth}
\begin{tabbing}
\quad\=\quad\=\quad\=\kill
\textbf{mientras} quede tiempo:\\
\> $s \leftarrow$ \textsc{InicialAleatoria}()\\
\> \textbf{repetir}\\
\>\> $s' \leftarrow \arg\min_{v\in N_1(s)\cup N_2(s)} f(v)$\\
\>\> \textbf{si} $f(s') < f(s)$: $s\leftarrow s'$\\
\> \textbf{hasta} que no haya mejora (óptimo local)\\
\> \textbf{si} $f(s)<f(s^*)$: $s^*\leftarrow s$\\
\textbf{retornar} $s^*$
\end{tabbing}
\column{0.47\textwidth}
\begin{itemize}
\item \textbf{Best Improvement}: se evalúan \emph{todos} los vecinos de $N_1$ y $N_2$ antes de mover.
\item \textbf{Evaluación incremental}: un movimiento solo cambia 2 turnos, así que\\
$\Delta=f(A')+f(B')-f(A)-f(B)$, con $f(A),f(B)$ guardados en caché.
\item Solo se aceptan mejoras estrictas.
\item Criterio de término: tiempo máximo o número de reinicios (parámetros).
\item Iteraciones informadas = movimientos aplicados en total.
\end{itemize}
\end{columns}
\end{frame}

%---------------------------------------------------------------
\begin{frame}{Movimiento N1: Mover un tramo}
\begin{columns}[T]
\column{0.5\textwidth}
\begin{tikzpicture}[x=0.4cm,y=0.8cm,font=\scriptsize]
  \node[left] at (0,2) {A}; \node[left] at (0,1) {B};
  \foreach \a/\b/\l in {1/3/a1,3.5/6/a2,6.5/8/j,10/12/a3}{\draw[fill=teal!35] (\a,1.8) rectangle (\b,2.2); \node at ({(\a+\b)/2},2) {\l};}
  \foreach \a/\b/\l in {2/4.5/b1,9/11/b2}{\draw[fill=acento2!35] (\a,0.8) rectangle (\b,1.2); \node at ({(\a+\b)/2},1) {\l};}
  \draw[-Stealth,thick,red] (7.25,1.75) to[bend right=20] (7.25,1.25);
  \draw[dashed,red] (6.5,0.8) rectangle (8,1.2);
  \node[right] at (0,0.1) {$|N_1| = n\cdot(m+1)$ (incluye un turno vacío nuevo)};
\end{tikzpicture}
\column{0.47\textwidth}
\textbf{¿Por qué?}
\begin{itemize}
\item \textbf{Conectividad}: con N1 se puede llegar desde cualquier asignación a cualquier otra (por ejemplo, pasando por ``un conductor por tramo''). Así se cumple que los movimientos permitan visitar todo el espacio de búsqueda.
\item Es el ajuste fino: reubica tramos sueltos para equilibrar turnos o vaciar uno corto.
\item Si el destino tiene traslape, el vecino no se genera.
\end{itemize}
\end{columns}
\end{frame}

%---------------------------------------------------------------
\begin{frame}{Movimiento N2: Cruce (intercambio de sufijos, 2-opt*)}
\begin{columns}[T]
\column{0.5\textwidth}
\begin{tikzpicture}[x=0.4cm,y=0.8cm,font=\scriptsize]
  \node[left] at (0,3) {A}; \node[left] at (0,2) {B};
  \foreach \a/\b in {1/3,3.5/5.5}{\draw[fill=teal!35] (\a,2.8) rectangle (\b,3.2);}
  \foreach \a/\b in {8/10,10.5/12}{\draw[fill=teal!15,dashed] (\a,2.8) rectangle (\b,3.2);}
  \foreach \a/\b in {2/4,4.5/6.5}{\draw[fill=acento2!35] (\a,1.8) rectangle (\b,2.2);}
  \foreach \a/\b in {7/9.5,13/14}{\draw[fill=acento2!15,dashed] (\a,1.8) rectangle (\b,2.2);}
  \draw[red,thick] (6.8,1.6) -- (6.8,3.4) node[above] {corte};
  \node[left] at (0,0.9) {A'}; \node[left] at (0,0.1) {B'};
  \foreach \a/\b in {1/3,3.5/5.5}{\draw[fill=teal!35] (\a,0.7) rectangle (\b,1.1);}
  \foreach \a/\b in {7/9.5,13/14}{\draw[fill=acento2!15,dashed] (\a,0.7) rectangle (\b,1.1);}
  \foreach \a/\b in {2/4,4.5/6.5}{\draw[fill=acento2!35] (\a,-0.1) rectangle (\b,0.3);}
  \foreach \a/\b in {8/10,10.5/12}{\draw[fill=teal!15,dashed] (\a,-0.1) rectangle (\b,0.3);}
\end{tikzpicture}

\small $A'=A[0..i)+B[k..]$, $B'=B[0..k)+A[i..]$. El corte $k$ se calcula a partir de $i$, así que $|N_2|=O(m^2 L)$.
\column{0.47\textwidth}
\textbf{¿Por qué?}
\begin{itemize}
\item Se basa en el intercambio de sufijos de De Leone et al.~[2] y en la recombinación de segmentos de Constantino et al.~[1].
\item Mueve \textbf{bloques completos} de trabajo en un solo paso, cosa que N1 necesitaría varios pasos para lograr, y varios de esos pasos empeoran el costo.
\item Si $B[0..k)$ queda vacío, B desaparece. Así se \textbf{fusionan turnos cortos}, que es lo que más ahorra porque cada turno cuesta al menos 432.
\item Mantiene juntos los encadenamientos de un mismo bus, lo que conserva la continuidad espacial.
\end{itemize}
\end{columns}
\end{frame}

%---------------------------------------------------------------
\begin{frame}[fragile]{Salida y resultados preliminares}
\begin{columns}[T]
\column{0.3\textwidth}
\begin{lstlisting}[language=bash]
9681.08 30.001 1 77343

1 12
2 12
3 12
...
130 6
\end{lstlisting}
\small \texttt{<FO> <tiempo> <semilla> <iter>}, luego \texttt{<tramo> <conductor>}. Se escribe en stdout y en \texttt{salidas/}.
\column{0.67\textwidth}
\scriptsize
\begin{tabular}{lrrrrrr}
\toprule
& \multicolumn{2}{c}{Inicial aleatoria} & \multicolumn{2}{c}{HC-BI + Restart (30 s)} & & \\
Instancia & cond. & FO & cond. & FO & mejora & reinicios\\
\midrule
V\_TCCC0130 & 27 & 13\,341,8 & 22 & 9\,681,1 & 27,4\% & 1178\\
V\_TCCC0251 & 49 & 24\,656,4 & 46 & 20\,228,7 & 18,0\% & 151\\
v\_cv412    & 81 & 40\,967,6 & 80 & 35\,318,9 & 13,8\% & 36\\
V\_TCCC0512 & 94 & 48\,385,5 & 93 & 40\,872,6 & 15,5\% & 17\\
V\_TCCC0761 & 134 & 70\,141,1 & 133 & 58\,596,0 & 16,5\% & 5\\
V\_TCCC1000 & 172 & 90\,485,6 & 172 & 76\,267,8 & 15,7\% & 3\\
\bottomrule
\end{tabular}

\vspace{0.3em}
\small Semilla 1. Todas las soluciones son factibles.
\begin{tikzpicture}
\begin{axis}[ybar,height=3.3cm,width=0.95\textwidth,bar width=7pt,
  symbolic x coords={130,251,412,512,761,1000},xtick=data,
  ylabel={\scriptsize mejora (\%)},xlabel={\scriptsize tramos},ymin=0,ymax=34,
  tick label style={font=\scriptsize},nodes near coords,
  every node near coord/.append style={font=\tiny},enlarge x limits=0.1]
\addplot[fill=acento!60] coordinates {(130,27.4) (251,18.0) (412,13.8) (512,15.5) (761,16.5) (1000,15.7)};
\end{axis}
\end{tikzpicture}
\end{columns}
\end{frame}

%---------------------------------------------------------------
\begin{frame}{Problemas encontrados y dudas}
\textbf{1. Meseta generada por el tiempo pagado mínimo.}
\begin{itemize}\small
\item Un turno con $W<432$ cuesta 432 sin importar cuántos tramos tenga. Quitarle tramos no mejora $f$ hasta que queda \emph{vacío}, y el HC-BI se estanca.
\item Por eso, en las instancias grandes casi no baja el número de conductores (172 $\to$ 172).
\item \textbf{Propuesta}: agregar un término secundario de desempate que premie vaciar los turnos cortos, por ejemplo $-\epsilon\sum_k W_k^2$, guiando la búsqueda por el tiempo ocioso como recomiendan Constantino et al.~[1].
\end{itemize}
\textbf{2. Costo de una iteración en instancias grandes.}
\begin{itemize}\small
\item Explorar $N_1\cup N_2$ completo cuesta $O(nmL + m^2L^2)$. En 2313 tramos, un solo descenso no termina en 20 s.
\item \textbf{Propuesta}: podar por ventana temporal, es decir, solo considerar turnos que se solapen en el tiempo con el tramo o el corte.
\end{itemize}
\textbf{Dudas para el ayudante}
\begin{itemize}\small
\item ¿La FO de la salida debe incluir la penalización por cambio de línea? Por ahora sí la incluye.
\item ¿Se admite más de un descanso no pagado por turno? Por ahora se supone uno, igual que en el informe.
\end{itemize}
\end{frame}

%---------------------------------------------------------------
\begin{frame}{Conclusiones y trabajo restante}
\textbf{Conclusiones}
\begin{itemize}
\item Ya funcionan la lectura, la evaluación, la solución inicial, la salida y una primera versión del HC-BI con reinicios, en todas las instancias.
\item La asignación directa con evaluación por turno en caché hace que cada vecino cueste $O(L)$.
\item El HC-BI mejora entre 14\% y 27\% sobre la solución inicial, pero la meseta de $P^{\min}$ limita la reducción de conductores.
\end{itemize}
\textbf{Trabajo restante (Entrega 2, 23 de noviembre)}
\begin{enumerate}
\item Término de desempate para escapar de la meseta de $P^{\min}$.
\item Poda de vecindarios para las instancias grandes (1253 a 2313 tramos).
\item Incorporar la retroalimentación de esta presentación.
\item Experimentos: varias semillas, comparación de $N_1$, $N_2$ y $N_1\cup N_2$, y comparación con los resultados de la literatura~[1].
\item Informe 2: representación, algoritmo, experimentos y resultados.
\end{enumerate}
\end{frame}

%---------------------------------------------------------------
\begin{frame}[fragile]{Uso de LLMs}
\small
\textbf{Herramienta:} Claude Code (Anthropic, modelo Claude Opus 5.5), como asistente de
programación bajo mi especificación y revisión.

\begin{itemize}
\item \textbf{Mis decisiones de diseño:} representación (asignación directa
      tramo$\to$conductor), inicialización aleatoria no-greedy, vecindarios N1 y N2,
      y \texttt{PESO\_VIOLACION}$=1000>432$ para que reparar una violación siempre
      convenga más que ahorrar un turno.
\item \textbf{Del asistente:} redacción del código C++ a partir de esa especificación,
      el \texttt{Makefile}, el \texttt{README} y borradores de estas diapositivas.
\item \textbf{Mi revisión:} ajuste del código, verificación a mano del pago nocturno
      (V\_TCCC0130, semilla~3, conductor~1), detección de la meseta de $P^{\min}$ y del
      costo por iteración en instancias grandes, y formulación de las dos mejoras
      propuestas.
\item \textbf{Validación:} ejemplo de 6 tramos (FO$=864$) y las 10 instancias.
\end{itemize}
\end{frame}

%---------------------------------------------------------------
\begin{frame}{Referencias}
\small
\begin{enumerate}
\item[{[1]}] A. A. Constantino, C. F. X. Mendonça, S. A. Araujo, D. Landa-Silva, R. Calvi, A. F. Santos. \emph{Solving a Large Real-world Bus Driver Scheduling Problem with a Multi-assignment based Heuristic Algorithm}. Journal of Universal Computer Science, 23(5):479--504, 2017.
\item[{[2]}] R. De Leone, P. Festa, E. Marchitto. \emph{A Bus Driver Scheduling Problem: a new mathematical model and a GRASP approximate solution}. Journal of Heuristics, 17:441--466, 2011.
\item[{[3]}] L. Kletzander, N. Musliu. \emph{Solving Large Real-Life Bus Driver Scheduling Problems with Complex Break Constraints}. ICAPS, 2020.
\end{enumerate}
\end{frame}

\end{document}
